%=======================================================================
%  2. CHALLENGE & SCOPING  (text from the Word report; audit edits only)
%=======================================================================
\section{Challenge \& Scoping}
\label{sec:scoping}

\subsection{Problem Definition}
\label{sec:problem}

Conventional cooling methods such as air conditioning worsen the Urban Heat Island
effect by displacing the warm air from inside the building to the surrounding area.
This local heat exchange causes the city to warm up more, increasing the need for
air conditioning, which ultimately results in a self-reinforcing feedback loop.

\subsection{Context and Constraints}
\label{sec:context}

All context values, benchmarks and engineering checks in this report refer to the
City of Zurich. Zürich was chosen for four reasons:

\begin{enumerate}
\item The city has adopted a heat-mitigation strategy, the Fachplanung
  Hitzeminderung~\cite{stadtzuerich2020}, with an implementation agenda, a guidance
  leaflet for building owners and thematic maps. The project therefore has a
  documented heat map and an albedo classification to refer to.
\item The city publishes an open climate analysis including geodata, which let us
  select a documented reference site.
\item MeteoSwiss operates long term measurement stations in and around the city.
\item Its building style and infrastructure are representative of the Swiss
  Mittelland, so results transfer to comparable Swiss cities.
\end{enumerate}

\subsubsection{Reference Building}
\label{sec:refbuilding}

For comparison and reference purposes a building inside of Zurich was chosen. The
chosen building is the Modellierungsgebiet 03 -- Geschlossene Randbebauung (closed
perimeter block). It consists of two outward-closed blocks of five stories including
a roof structure and an enclosed private courtyard. It is located in a dense
district close to the city center. The city's assessment is that the reduced air
exchange and lack of cold air supply makes this building overheat comparatively
quickly~\cite[pp.~148--151]{stadtzuerich2020}.

\subsubsection{Constraints}
\label{sec:constraints}

Table~\ref{tab:constraints} separates hard constraints where a violation
disqualifies a concept, from soft constraints which weight the evaluation but do not
necessarily disqualify a concept.

\begin{table}[!htb]
\centering
\caption{Context conditions and constraints.}
\label{tab:constraints}
\renewcommand{\arraystretch}{1.15}
\small
\begin{tabularx}{\textwidth}{L{2.7cm} L{0.9cm} X}
\toprule
\textbf{Type} & \textbf{Class} & \textbf{Constraint} \\
\midrule
Energy & hard & Limit auxiliary electrical energy for the cooling function. \\
Water & hard & No unnecessary water waste. \\
Townscape & hard & The result must be acceptable in the Zurich townscape. Buildings may be listed; heritage protection, building law and neighbor rights can restrict visible changes~\cite[p.~128]{stadtzuerich2020}. \\
Glare & hard & No specular glare towards streets, neighboring flats or traffic, which can impede day to day life. \\
Climate & hard & Must survive the Swiss annual cycle: frost to about \qty{-10}{\celsius}, freeze--thaw cycles, driving rain, UV-Light, and periods of several dry weeks in summer (team design assumptions). \\
Winter & hard & Must not increase heating demand. A high solar reflectance also reduces useful winter solar gains; this trade-off has to be quantified. \\
Existing building & soft & Preferable, but not necessary. Concepts can also be developed for only new buildings. \\
Installation & soft & Must be installable by ordinary trades on an occupied building, ideally without full scaffolding. \\
Maintenance & soft & Must remain functional with little maintenance. \\
Solar energy & soft & Should not unnecessarily compromise roof- or façade-integrated photovoltaics. \\
Structural & soft & Limited additional dead load and build-up thickness on an existing envelope. \\
Embodied emissions & soft & The greenhouse gas embodied in the added infrastructure must be repaid by the avoided cooling demand within the service life, measured in kg CO\textsubscript{2}-eq per m\textsuperscript{2}. \\
Economic & soft & Cost per square meter should be in the range of a normal refurbishment. \\
\bottomrule
\end{tabularx}
\end{table}

\subsection{System View and Project Boundary}
\label{sec:systemview}

\subsubsection{Super-system}

The chosen super system is the city and its existing infrastructure.

\subsubsection{System of interest}

The core interest of this project focuses on heat-regulation of the outer layer of
a building including the façade and the roof.

\subsubsection{Parallel systems}

\begin{itemize}
\item Stormwater and sewer network
\item Street and pavement network
\item Power grid and building energy system
\item Residents
\item Cars, Bicycles, Public transport and other forms of transportation
\end{itemize}

\subsubsection{Subsystems}

\begin{itemize}
\item Outer surface
\item Storage and insulation layer
\item Water paths
\item Attachment to the substrate
\end{itemize}

Table~\ref{tab:inout} states which aspects are scoped in or out.

\begin{table}[!htb]
\centering
\caption{Project boundary: what is deliberately in and out of the project.}
\label{tab:inout}
\renewcommand{\arraystretch}{1.15}
\small
\begin{tabularx}{\textwidth}{X X}
\toprule
\textbf{IN the project} & \textbf{OUT of the project} \\
\midrule
Prevent heat buildup and remove heat from the building façade. & Redesign of the building itself. Reuse of the stored energy. Internal cooling. \\
\bottomrule
\end{tabularx}
\end{table}

\subsection{Design Question}
\label{sec:designquestion}

\subsubsection{Candidate questions considered}

Table~\ref{tab:questions} shows the rejected candidates.

\begin{table}[!htb]
\centering
\caption{Design questions considered, and why they were rejected.}
\label{tab:questions}
\renewcommand{\arraystretch}{1.15}
\small
\begin{tabularx}{\textwidth}{L{0.5cm} L{5.4cm} X}
\toprule
\textbf{\#} & \textbf{Candidate} & \textbf{Assessment} \\
\midrule
1 & How do we prevent heat build-up in urban areas? & Too broad. No place, no time, no intervention point. The equivalent of \enquote{How might we end world hunger?} \\
2 & How could we use the energy or heat build-up positively while lowering temperature? & Right idea, but unanchored: no location, and \enquote{positively} does not say for what. \\
3 & How might we get the heat that Zurich's existing buildings store during the day back out of the city before the next morning, without using energy? & Well-formed and highly testable, but it prescribes when and how the heat has to leave and leaves little room for variation. Its intervention point, the existing buildings, is taken over into the selected question. \\
4 & How can we prevent heat waves causing heat islands? & Factually wrong. Heat waves do not cause heat islands. The urban heat island is a permanent consequence of urban infrastructure and exists on ordinary summer days. A heat wave makes it dangerous, but it does not create it. \\
\bottomrule
\end{tabularx}
\end{table}

\subsubsection{Selected design question}

\emph{\textbf{How might we reduce the urban heat island effect in cities like
Zurich by changing the façade and roof of existing buildings?}}

The added clause fixes the intervention point at the system of interest of
Section~\ref{sec:systemview}. It allows some freedom in how to reduce the urban heat
island effect, and it takes over the mechanism of candidate 3 without prescribing when
and how the heat has to leave.

\needspace{16\baselineskip}
\subsection{Functions and Functional Diagram}
\label{sec:functions}

\subsubsection{Functions}

Table~\ref{tab:functions} displays the chosen functions and their respective
taxonomy link.

\begin{table}[!htb]
\centering
\caption{Functions list.}
\label{tab:functions}
\renewcommand{\arraystretch}{1.15}
\small
\begin{tabularx}{\textwidth}{L{0.8cm} L{4.2cm} X}
\toprule
\textbf{ID} & \textbf{Function (verb)} & \textbf{Taxonomy link -- Explanation} \\
\midrule
F1 & Reflect short-wave radiation & Prevent solar radiation -- Avoid absorbing solar energy (Visible \& near infrared light). \\
F2 & Shade a surface & Prevent excess temperature -- Reduce direct light exposure on structural surfaces. \\
F3 & Dissipate / Transfer heat & Modify thermal conditions -- Transfer thermal energy via transmission or thermal radiation. \\
F4 & Evaporate liquid & Regulate physical processes -- Heat extraction using moisture or water \\
F5 & Resist degradation & Prevent structural breakdown -- Maintain long-term durability against UV, rain, and frost. \\
\bottomrule
\end{tabularx}
\end{table}

\subsubsection{Functional Diagram}
\label{sec:funcdiagram}

Figure~\ref{fig:funcdiagram} draws the functions of Table~\ref{tab:functions} as a
flow picture. Incoming shortwave radiation is either reflected (F1) or the surface is
shaded (F2); what is absorbed raises the surface temperature and leaves again by F3,
to the sky and to the street-canyon air, or by F4, where liquid water enters the
layer and leaves as vapour. Whatever neither path removes flows into the building.
The surface temperature is the signal that sets how strongly F3 and F4 act.

% Functional diagram: orthogonal (straight-line) layout on a 4-column grid.
% Columns: inputs (x=0), functions F1/F2 + state (x=3.4), functions F3/F4 (x=6.8), sinks (x=10.4).
\begin{figure}[htbp]
\centering
\scalebox{0.9}{%
\begin{tikzpicture}[x=1cm,y=1cm,
  lbl/.style={font=\scriptsize,fill=white,inner sep=1.5pt},
  dot/.style={circle,fill=black!65,inner sep=1.3pt}]

% ---------------------------------------------------------------- nodes
\node[funcbox,fill=yellow!25,text width=1.9cm,minimum height=9mm] (sun) at (0,2.6)
  {Solar\\ irradiance};
\node[funcbox,text width=2.4cm,minimum height=11mm] (f12) at (3.4,3.7)
  {\textbf{F1} Reflect short-wave radiation\\[1pt] \textbf{F2} Shade a surface};
\node[funcbox,fill=red!8,text width=2.4cm,minimum height=11mm] (ts) at (3.4,1.4)
  {Absorbed heat\\[1pt] \itshape state: surface temperature $T_s$};
\node[funcbox,text width=2.4cm,minimum height=9mm] (f3) at (6.8,3.7)
  {\textbf{F3} Dissipate / transfer heat};
\node[funcbox,text width=2.4cm,minimum height=9mm] (f4) at (6.8,0.5)
  {\textbf{F4} Evaporate liquid};
\node[funcbox,fill=zhawblue!10,text width=1.9cm,minimum height=9mm] (wat) at (6.8,-1.1)
  {Water in};
\node[funcbox,fill=gray!12,text width=2.0cm,minimum height=9mm] (sky) at (10.4,3.7) {Sky};
\node[funcbox,fill=gray!12,text width=2.0cm,minimum height=9mm] (air) at (10.4,2.3) {Street-canyon air};
\node[funcbox,fill=gray!12,text width=2.0cm,minimum height=9mm] (vap) at (10.4,0.5) {Vapour to air};
\node[funcbox,fill=gray!12,text width=2.0cm,minimum height=9mm] (bld) at (10.4,-2.1) {Building interior};

% ---------------------------------------------------------- energy flows
% sun -> F1/F2 and -> absorbed heat (branch b1)
\coordinate (b1) at (1.7,2.6);
\draw[flowarr] (sun.east) -- (b1) |- (f12.west);
\draw[flowarr] (b1) |- (ts.west);
\node[dot] at (b1) {};
\node[lbl,anchor=east] at (1.58,2.0) {absorbed};

% reflected: over the top, straight down into the sky
\draw[flowarr] (f12.north) -- ++(0,1.1) -| (sky.north);
\node[lbl] at (6.9,5.35) {reflected};

% absorbed heat -> F3 and F4 (branch b2)
\coordinate (b2) at (5.0,1.4);
\draw[flowarr] (ts.east) -- (b2) |- (f3.west);
\draw[flowarr] (b2) |- (f4.west);
\node[dot] at (b2) {};

% F3 -> sky (radiation) and -> canyon air (convection), branch b3
\coordinate (b3) at (8.6,3.7);
\draw[flowarr] (f3.east) -- (b3) -- (sky.west);
\draw[flowarr] (b3) |- (air.west);
\node[dot] at (b3) {};
\node[lbl] at (8.7,4.0) {radiation};
\node[lbl,rotate=90] at (8.6,3.0) {convection};

% residual: absorbed heat -> building interior (to be minimised)
\draw[flowarr,draw=red!70] (ts.south) |- (bld.west);
\node[lbl,text=red!70] at (6.8,-1.82) {residual heat into the building};

% -------------------------------------------------------- material flows
\draw[wateharr] (wat.north) -- (f4.south);
\draw[wateharr] (f4.east) -- (vap.west);

% ------------------------------------------------------ information flow
% T_s drives F3 and F4: orange dashed lines run parallel to the energy paths
\draw[infoarr,transform canvas={xshift=3.5pt,yshift=-3.5pt}] (ts.east) -- (b2) |- (f3.west);
\draw[infoarr,transform canvas={xshift=3.5pt,yshift=-3.5pt}] (b2) |- (f4.west);
\node[lbl,text=orange!85!black,font=\scriptsize\bfseries] at (5.85,2.55) {$T_s$ signal};

% ------------------------------------------------------- system boundary
\begin{scope}[on background layer]
  \node[draw=zhawblue,dashed,rounded corners=4pt,fit=(f12)(ts)(f3)(f4),
        inner sep=9pt] (fb) {};
\end{scope}
\node[font=\scriptsize\bfseries,color=zhawblue,anchor=east]
  at ($(fb.north east)+(-0.1,0.28)$) {SYSTEM BOUNDARY};

% ---------------------------------------------------------------- legend
\node[font=\scriptsize,anchor=west,align=left] at (-1.1,-3.1)
  {\textcolor{black!65}{\textbf{---}} energy flow \quad
   \textcolor{zhawblue}{\textbf{--\,--}} material flow (water, vapour) \quad
   \textcolor{orange!85!black}{\textbf{--\,--}} information flow \quad
   \textcolor{red!70}{\textbf{---}} flow to be minimised};
\end{tikzpicture}}
\caption{Solution-neutral functional diagram; the system boundary is the outer layer of the building. Solid arrows are energy flows, blue
dashed arrows the material flow of F4 (water in, vapour out), the orange dashed
arrows the information flow: the surface temperature drives the strength of F3 and
F4. F5 (resist degradation) acts on everything inside the boundary and is not
drawn.}
\label{fig:funcdiagram}
\end{figure}


\subsection{Biologized Questions}
\label{sec:bioq}

\begin{enumerate}
\item How does nature reflect incoming shortwave radiation?
\item How does nature shade a surface even in the desert?
\item How does nature dissipate excess heat?
\item How does nature transfer excess heat?
\item How does nature dissipate excess heat using moisture?
\item How does nature keep a functional surface clean without cleaning it?
\end{enumerate}

\subsection{Performance Criteria and KPIs}
\label{sec:kpi}

\subsubsection{KPIs}

The Albedo index defines a value between 0 -- 1 where 0 is total absorption of all
incoming radiation and 1 is the total reflection of all incoming
radiation~\cite{nasa_albedo}. Shading provides an alternative way to keep a surface
cool~\cite{dussadee2018}. Water evaporation is an effective solution to remove heat
using moisture as means of heat transfer; the aim is to reduce the needed amount of
water evaporation to conserve water.

Table~\ref{tab:kpi} defines one KPI per function, with the unit in which we compare
concepts, the baseline of the untreated reference block, the target, and the
measurement we would use to check it. The baseline class and the target of KPI1 come
from the albedo classification of the Fachplanung, which rates 0.30 and below as low
and 0.80 and above as high~\cite[p.~128]{stadtzuerich2020}. The target of KPI5
follows from the Water constraint in Table~\ref{tab:constraints}. The targets of
KPI2, KPI4 and KPI6 are team working targets: we set them to have a number to test
against, and they are not city values. KPI3 and KPI4 both belong to F3 but measure
different things: KPI3 counts the heat that actually leaves the façade through a
removal path, which is the quantity the engineering checks in
Section~\ref{sec:check1} and Section~\ref{sec:check2} estimate; KPI4 is the effect on
the surface that a pedestrian is exposed to.

\begin{table}[H]
\centering
\caption{KPIs. Targets marked \enquote{team} were set by the team; the KPI1
target and baseline class come from the Fachplanung~\cite[p.~128]{stadtzuerich2020}.}
\label{tab:kpi}
\renewcommand{\arraystretch}{1.0}
\scriptsize
\begin{tabularx}{\textwidth}{L{0.7cm} L{0.45cm} L{3.3cm} L{1.5cm} L{2.3cm} L{2.1cm} X}
\toprule
\textbf{\#} & \textbf{Fn} & \textbf{KPI} & \textbf{Unit} & \textbf{Baseline} &
\textbf{Target} & \textbf{Measurement} \\
\midrule
KPI1 & F1 & Solar reflectance of the exposed surface & -- (0--1) &
  Existing render and roof of the reference block, class \enquote{low}
  ($\leq 0.30$) &
  $\geq 0.80$, the city's \enquote{high} class &
  Spectrophotometer or pyranometer pair, new and after weathering \\
KPI2 & F2 & Share of the south and west façade area shaded between 12:00 and
  18:00 on the design day (21~June, clear sky) & \unit{\percent} &
  \qty{0}{\percent} (bare façade) &
  $\geq \qty{70}{\percent}$ (team) &
  Shadow simulation on the 3D model, then site photos \\
KPI3 & F3 & Heat removed from the façade and released outside the street canyon &
  \unit{\kilo\watt} total; \unit{\watt\per\square\metre} of connected façade &
  0 (no removal path today) &
  $\geq$ the collected façade load, taken as \qty{100}{\watt\per\square\metre}
  in Section~\ref{sec:check1} &
  Calorimetric: flow $\times$ $c_p$ $\times$ $\Delta T$ on the loop; for a
  coating, surface heat-flux plates \\
KPI4 & F3 & Reduction of the façade surface temperature at 18:00 on the design
  day, against the untreated reference façade & \unit{\kelvin} &
  \qty{0}{\kelvin} &
  $\geq \qty{5}{\kelvin}$ (team) &
  Infrared thermometer on treated and untreated bays of the same façade \\
KPI5 & F4 & Water consumed in operation & \unit{\litre} per \unit{\square\metre}
  façade per day &
  0 (the reference façade uses none) &
  0 (Water constraint, Table~\ref{tab:constraints}) &
  Water meter on the make-up line \\
KPI6 & F5 & Performance retained after 20 annual cycles, for KPI1 and KPI3 &
  \unit{\percent} of initial & -- &
  $\geq \qty{90}{\percent}$ (team) &
  Accelerated weathering (UV, freeze--thaw, driving rain), then repeat KPI1
  and KPI3 \\
\bottomrule
\end{tabularx}
\end{table}


\subsubsection{Life Principles}
\label{sec:lp}

The life principles which should reflect the end solution include:

\begin{itemize}
\item Be resource efficient: minimize the use of external resources
\item Adapt to changing conditions: adapt to day/night time as well as season cycles
\item Use life-friendly chemistry: refrain the use of toxic chemicals or materials
\end{itemize}
